\section{Signature specialisation}
\label{sec:specialization}

After monomorphisation with the \emph{drop} policy, the guard
structure of the problem is ground: every $\hastype$-atom has a
closed type encoding as its second argument.  The problem is
nevertheless still written in the ``apply-everything'' dialect: a
single binary $\app$, a single binary $\hastype$, a single unary
$\PR$.  \emph{Signature specialisation} rewrites it, occurrence by
occurrence, into an ordinary first-order signature: unary guard
predicates $\guard{\tau}$, $n$-ary function symbols, $n$-ary
predicate symbols.  Two constraints shape the definition.  First,
the translation is to remain \emph{shallow}: specialisation is a
per-occurrence renaming and introduces \emph{no bridging axioms}
connecting the specialised symbols back to $\hastype$, $\app$ or
$\PR$.\footnote{The CoqHammer implementation currently prints
$\hastype(x, c\,t_1 \cdots t_n)$ as $c_{\$t}(t_1, \ldots, t_n, x)$
\emph{and} emits an unconditional bridging axiom $\forall \tup{x}
y.\ c_{\$t}(\tup{x}, y) \leftrightarrow \hastype(y, c\,\tup{x})$
\cite[\S5.5]{CzajkaKaliszyk2018}, which keeps the generic apparatus
alive in every problem.  The design here eliminates the bridge; the
price is a completeness loss quantified by
\cref{thm:spec-faithful}: none on the first-order core, only outside
it.}  Second, soundness must be unconditional, while
proof-preservation may legitimately be restricted to the fragment
where the rewriting is total.

Throughout this section $(\Phi, \gamma)$ is a translated problem
(\cref{def:transl-problem}), possibly after monomorphisation
(\cref{def:mono-op}).

\begin{definition}[Specialisation maps]
\label{def:spec}
Let $\mathcal{T}$ be a finite set of closed $\CICm$ types and
kind-level terms (in practice: all closed second arguments of
$\hastype$-atoms in $\Phi \cup \{\gamma\}$, plus $\SProp$), and let
$\mathcal{H}$ be a finite set of pairs $(c, n)$ of a constant
(original or lifted) and an arity.  The specialisation
$\spec{}$ of a formula is the result of exhaustively applying the
following rewrites, innermost first, to its subterms and atoms:
\begin{enumerate}[label=(S\arabic*),leftmargin=3em]
\item \emph{Type-argument fusion:} an occurrence $\hat{c} \app
  \trCn{\tau_1} \app \cdots \app \trCn{\tau_j} \app t_1 \app \cdots$
  where $c$'s type begins with $j$ binders of domain $\STypeI$ and
  all $\tau_i \in \mathcal{T}$ are closed, is rewritten to
  $\hat{c}_{\langle\tup{\tau}\rangle} \app t_1 \app \cdots$, where
  $\hat{c}_{\langle\tup{\tau}\rangle}$ is a fresh constant indexed
  by $(c, \tup{\tau})$.
\item \emph{Guard specialisation:} an atom $\hastype(t, u)$ where
  $u$ is syntactically $\trCn{\tau}$ (after (S1) on $u$'s
  arguments, if any) for some $\tau \in \mathcal{T}$, is rewritten
  to $\guard{\tau}(t)$, a fresh unary predicate indexed by $\tau$.
\item \emph{Arity flattening:} an occurrence $d \app t_1 \app \cdots
  \app t_n$ with head a constant $d$ (possibly produced by (S1)) and
  $(d, n) \in \mathcal{H}$, not itself an argument of a further
  $\app$ with the same head chain, is rewritten to $d^{(n)}(t_1,
  \ldots, t_n)$, a fresh $n$-ary function symbol; for $n = 0$ the
  constant is kept as is.
\item \emph{Predicate flattening:} an atom $\PR(d^{(n)}(t_1, \ldots,
  t_n))$ or $\PR(d)$ with $d$ a constant is rewritten to
  $d_{\mathsf{P}}^{(n)}(t_1, \ldots, t_n)$, a fresh $n$-ary
  predicate symbol.
\end{enumerate}
$\spec{}(\Phi, \gamma) = (\setof{\spec{}(\varphi)}{\varphi \in
\Phi}, \spec{}(\gamma))$.  The rewriting is \emph{total} on
$(\Phi, \gamma)$ if in the result the symbols $\hastype$, $\PR$ and
$\app$ do not occur at all.
\end{definition}

(S3) and (S4) are the classical optimisations of Meng and
Paulson~\cite{MengPaulson2008}, already performed by CoqHammer
\cite[\S5.5]{CzajkaKaliszyk2018}; (S1) is the symbol-discrimination
step of monomorphisation in the sense of
\cite[\S5.6]{Blanchette2016} and \textsf{Dis}
of~\cite{BobotPaskevich2011}; (S2) is the guard-specialisation this
paper is after.  All four are instances of one soundness principle:

\begin{theorem}[Soundness of specialisation]
\label{thm:spec-sound}
Let $(\Phi, \gamma)$ be $\MM$-sound via the expansion $\FM$ of
$\FMstar$ (\cref{def:si}).  Then $\spec{}(\Phi, \gamma)$ is
$\MM$-sound.
\end{theorem}

\begin{proof}
Expand $\FM$ to $\FM'$ by interpreting the fresh symbols by the
value functions of the patterns they replace:
\begin{align*}
  \hat{c}_{\langle\tup{\tau}\rangle} &\mapsto
    \ap(\cdots \ap(\sem{c}, \sem{\tau_1}) \cdots, \sem{\tau_j}), \\
  \guard{\tau} &\mapsto \setof{a \in \Dom}{a \in \sem{\tau}}, \\
  d^{(n)} &\mapsto (a_1, \ldots, a_n) \mapsto
    \ap(\cdots \ap(d^{\FM}, a_1)\cdots, a_n), \\
  d^{(n)}_{\mathsf{P}} &\mapsto
    \setof{(a_1, \ldots, a_n)}{\emptyset \in
      \ap(\cdots \ap(d^{\FM}, a_1)\cdots, a_n)} .
\end{align*}
By construction, each rewrite step replaces a term (resp.\ atom) by
one with the same value (resp.\ truth value) in $\FM'$ under every
valuation: for (S1) and (S3) this is the definition of the fresh
symbol's interpretation; for (S2), $\hastype^{\FM}(a, \sem{\tau})$
iff $a \in \sem{\tau}$ iff $\guard{\tau}^{\FM'}(a)$, using that
$\trCn{\tau}$ evaluates to $\sem{\tau}$ (\cref{lem:compositionality},
extended to (S1)-images by the first line); for (S4), $\PR^{\FM}(b)$
iff $\emptyset \in b$.  By induction on formulas, $\FM' \folmodels
\varphi \leftrightarrow \spec{}(\varphi)$ pointwise for every
$\varphi \in \Phi \cup \{\gamma\}$.  Hence $\FM' \folmodels
\spec{}(\Phi)$, and $\FM' \folmodels \spec{}(\gamma)$ implies $\FM
\folmodels \gamma$, which implies validity of the pending goal.
\end{proof}

Soundness costs nothing; the substantive question is whether
specialisation \emph{loses} proofs, given that no bridging axioms
tie $\guard{\tau}$ to $\hastype$.  On the first-order core it does
not:

\begin{lemma}[Totality on the first-order core]
\label{lem:spec-total}
Suppose every premise of the (drop-policy monomorphised) problem is
monomorphic, and every quantified variable in the premises and goal
has as its type a closed data type that is not a product, with no
quantification over propositions.  Before specialising,
\emph{prune} the axiom store: remove the universe atom (D4) (needed
only to discharge $\hastype(\cdot, \lift{\STypeI})$ guards, which
do not occur in monomorphic problems) and the application-typing
axioms ($\Ax$-app), whose bodies contain variable-headed
applications $z \app x$; these axioms exist to let a
\emph{quantified} function variable be applied, and product-typed
quantified variables are excluded here, while for constant-headed
terms their derivational content is subsumed by the
$\mathrm{Flat}$ axioms.  Pruning removes premises, so it trivially
preserves $\MM$-soundness.  Take $\mathcal{T}$ to be all closed
second arguments of surviving $\hastype$-atoms (including
membership atoms at product types, whose encodings are closed) and
$\mathcal{H}$ all (head, arity) pairs of constant-headed
application occurrences.  Then $\spec{}$ is total on the pruned
translated problem.
\end{lemma}

\begin{proof}
We check that no occurrence of $\hastype$, $\PR$, $\app$ survives.
$\hastype$-atoms: guards of quantified variables have closed
non-product data types, whose encodings are closed terms in
$\mathcal{T}$, so (S2) applies; the surviving typing atoms and the
$\mathrm{Flat}$ hypotheses and conclusions have closed non-product
type encodings (monomorphic declarations after instantiating the
implicit premises, \cref{def:mono-op}), so (S2) applies.  $\app$:
after pruning, every application occurrence has a constant head ---
variable heads arise only from variables of product or $\SProp$
type, excluded in the premises and goal, and from the ($\Ax$-app)
axioms, which are pruned --- so (S1)/(S3) apply exhaustively.
$\PR$: after (S3), every $\PR$-argument is a constant-headed
flattened term, so (S4) applies.
\end{proof}

\begin{theorem}[Faithfulness of total specialisation]
\label{thm:spec-faithful}
If $\spec{}$ is total on $(\Phi, \gamma)$ and $\Phi$ contains no
$\approx$-equations between type-level terms (guaranteed for
translated problems by \cref{lem:no-type-eq}), then
\[
  \Phi \folproves \gamma
  \qquad\Longleftrightarrow\qquad
  \spec{}(\Phi) \folproves \spec{}(\gamma) .
\]
\end{theorem}

\begin{proof}
By completeness it suffices to transfer models in both directions.

($\Leftarrow$ direction of provability, i.e.\ model transfer from
$\Phi$-side to $\spec{}$-side)  Let $\FM \folmodels \Phi \cup
\{\neg\gamma\}$.  Interpret the fresh symbols over the same domain
by the value functions of their patterns, exactly as in the proof of
\cref{thm:spec-sound} but with $\FM$ in place of $\FMstar$'s
expansion.  The same pointwise equivalence gives a model of
$\spec{}(\Phi) \cup \{\neg\spec{}(\gamma)\}$.

($\Rightarrow$ direction, model transfer from $\spec{}$-side)  Let
$\FM' \folmodels \spec{}(\Phi) \cup \{\neg\spec{}(\gamma)\}$ with
domain $\Dom'$.  By totality, the symbols $\hastype$, $\PR$, $\app$,
the raw constants with nonzero occurring arity, and the type-encoding
constants do not occur in the specialised problem, so we may
interpret them freely; $0$-ary raw constants that survive
specialisation unchanged keep their $\FM'$-interpretation.  Build
$\FM$ with domain
\[
  \Dom \;\coloneqq\; \Dom' \;\uplus\;
  \setof{\mathbf{t}_\tau}{\tau \in \mathcal{T}^{\downarrow}}
  \;\uplus\;
  \setof{\mathbf{p}_{d, \tup{a}}}{(d,n) \in \mathcal{H},\
    \tup{a} \in (\Dom')^{<n}}
\]
where $\mathcal{T}^{\downarrow}$ is the subterm closure of
$\mathcal{T}$ (finitely many fresh points $\mathbf{t}_\tau$, one per
closed type-level term, pairwise distinct and distinct from
$\Dom'$), and the $\mathbf{p}$-points represent partial
applications.  Interpret: each type-encoding constant and each
application step among type encodings so that the value of
$\trCn{\tau}$ is $\mathbf{t}_\tau$ (possible since the
$\mathbf{t}$-points are free and the syntax is finite: interpret
$\hat{I}$ as $\mathbf{t}_I$, $\ap(\mathbf{t}_{I\,\tup{\sigma}^{<}},
\mathbf{t}_\sigma) = \mathbf{t}_{I\,\tup{\sigma}^{<}\sigma}$, etc.,
with an arbitrary default elsewhere); $\ap$ on a raw constant head
$d$ and arguments $\tup{a} \in \Dom'$ climbs the
$\mathbf{p}$-points and, at the $n$-th argument for $(d, n) \in
\mathcal{H}$, returns $(d^{(n)})^{\FM'}(\tup{a})$ if $d^{(n)}$
occurs in the specialised problem, and otherwise a fresh terminal
$\mathbf{p}$-point; fused heads
$\hat{c}_{\langle\tup{\tau}\rangle}$ are reached as the value of
$\hat{c}$ applied to $\mathbf{t}_{\tau_1}, \ldots$, again via
dedicated $\mathbf{p}$-points;
\begin{align*}
  \hastype^{\FM}(a, b) \;&:\iff\;
    b = \mathbf{t}_\tau \text{ for some } \tau \in \mathcal{T}
    \text{ and } a \in \guard{\tau}^{\FM'},\\
  \PR^{\FM}(a) \;&:\iff\;
    a = \text{terminal point of } (d, \tup{a}) \text{ and }
    (d^{(n)}_{\mathsf{P}})^{\FM'}(\tup{a}),
\end{align*}
with $\PR$ false elsewhere and $\guard{\tau}^{\FM'} \subseteq
\Dom'$ read as a subset of $\Dom$.

We claim $\FM \folmodels \varphi \iff \FM' \folmodels
\spec{}(\varphi)$ for every $\varphi \in \Phi \cup \{\gamma\}$ and,
more generally, for every subformula under valuations $\nu$ sending
each free variable into $\Dom'$.  The proof is by induction on
$\varphi$, with an auxiliary term-level statement: for every term
$t$ occurring in $\varphi$ that is (i) an argument of $\approx$, of
a specialised guard atom, or of a flattened symbol, the
$\FM$-value of $t$ under $\nu$ equals the $\FM'$-value of
$\spec{}(t)$ under $\nu$.  For the term statement, note that such
$t$ is a fully applied constant-headed chain or a variable (by
totality), and the interpretation of $\ap$ was constructed to make
the chain land on exactly $(d^{(n)})^{\FM'}(\tup{a})$; variables are
shared.  Atoms: $\approx$-atoms compare $\Dom'$-values equal in both
structures --- here the absence of type-level equations matters:
both sides are data-level chains landing in $\Dom'$, never
$\mathbf{t}$- or $\mathbf{p}$-points; guard atoms $\hastype(t,
\trCn{\tau})$ hold in $\FM$ iff $\guard{\tau}^{\FM'}$ holds of the
(shared) value of $t$, which is the truth condition of the
specialised atom; $\PR$-atoms likewise via the terminal-point
clause.  Quantifiers: every quantifier of $\varphi$ is guarded
(inspection of \cref{def:translation}: (F2), (F5), the guarded
closures of all axiom schemes) by a guard whose $\FM$-extension is
$\guard{\tau}^{\FM'} \subseteq \Dom'$; hence the effective range of
every quantifier in $\FM$ is inside $\Dom'$, the valuations stay in
$\Dom'$, and the induction hypothesis applies.  (Universally
quantified formulas hold vacuously at the fresh points, which fail
every guard; existentials find their witnesses in $\Dom'$ on both
sides.)

Thus $\FM \folmodels \Phi \cup \{\neg\gamma\}$, completing the
transfer.
\end{proof}

\begin{remark}[Outside the first-order core]
\label{rem:spec-partial}
When higher-order quantification remains (e.g.\ a monomorphic
premise $\forall f{:}\mathsf{nat} \to \mathsf{nat}.\,\ldots$),
occurrences $f \app x$ with variable head survive specialisation,
$\app$ stays in the signature, and \cref{thm:spec-faithful} does
not apply --- specialisation remains sound
(\cref{thm:spec-sound}) but may lose proofs that exploited the
identity of $\hastype(\cdot, \lift{\mathsf{nat} \to \mathsf{nat}})$
with the guard of some other occurrence.  The bridge axioms of the
implementation buy back exactly these proofs at the cost of
reintroducing the generic apparatus; under the shallow-translation
constraint we accept the loss and note that on hammer benchmarks
the essentially first-order goals dominate
\cite{CzajkaKaliszyk2018}.
\end{remark}

\begin{remark}[Open guard shapes]
\label{rem:spec-shapes}
(S2) requires the guard's type argument to be closed.  A guard
$\hastype(y, \trCn{\mathsf{list}\,(Q\,x)})$ with a free data
variable $x$ (arising under a binder) can be specialised too, by a
\emph{shape} predicate: rewrite to
$\guard{\mathsf{list}\,(Q\,\cdot)}(x, y)$, a fresh binary predicate
indexed by the type pattern with its free variables abstracted.
\Cref{thm:spec-sound} extends verbatim --- interpret the shape
predicate by the value function $(a, b) \mapsto (b \in
\sem{\mathsf{list}\,(Q\,x)}_{[x \mapsto a]})$ --- so shape
specialisation is sound; we do not claim faithfulness for it.  This
settles, affirmatively for soundness, the question raised in the
CoqHammer development notes of whether guards over open types may
be compressed into parameterised predicates.
\end{remark}
